%% kapa: exam=1 round=19 year=2008 period=1 form=A subject=civil_law subject_ko=민법 seq=4 number=4
%% source: Q-Net 감정평가사 기출문제 (cst00309) artlSeq=1001962 file=extracted/1008/1001962/2008년 제19회 감정평가사 1차시험 기출문제/1교시.hwp
%% conversion: kordoc 4.24.51 -> tools/gen1.py (auto); review=proofread-sonnet; flags=none
\begin{question}{4}
민법상 법인의 불법행위능력에 관한 설명으로 틀린 것은? (다툼이 있는 경우에는 판례에 의함)
\begin{choices}
  \item 사단법인에서 특별대리인이나 직무대행자는 법인의 대표기관에 해당하지만 사원총회나 감사는 대표기관이 아니다.
  \item 법인의 불법행위책임을 야기한 이사 기타 대표자는 자기의 손해배상책임을 면하지 못한다.
  \item 대표기관의 직무와 밀접한 관련이 있을 뿐만 아니라 외형상으로도 대표기관의 직무범위 내의 행위로 볼 수 있으면 직무관련성이 인정된다.
  \item 대표기관이 개인적 이익을 위하여 대표권한을 남용한 경우도 직무에 관한 행위에 해당하지만, 상대방이 그 사정을 알았던 경우에는 법인이 상대방의 악의를 증명하여 그 효과를 부인할 수 있다.
  \item 대표기관의 행위가 직무집행에 관한 것이 아니라는 점에 대하여 상대방이 선의이면 중대한 과실이 있더라도 법인의 불법행위책임이 인정된다.
\end{choices}
\end{question}
\answer{5}
